\chapter{What Poplog is, and why it is unusual}

Poplog is not a language. It is a \emph{language toolkit} that happens to
ship with five languages already built on it: Pop-11, Prolog, Common Lisp,
Standard ML, and --- new in this fork --- Forth. All five share one virtual
machine, one garbage collector, one heap, and one process. You can start in
Prolog, drop into Pop-11 to write a predicate, and hand the result back,
without leaving the image or crossing a foreign-function boundary.

That design is forty years old and still rare. Most systems that host several
languages do it by compiling them all to a bytecode interpreter. Poplog does
it by compiling them all, incrementally, to \textbf{native machine code}.

\section{The two-level virtual machine}

Poplog's central idea is a split that lets language implementation and
machine porting proceed independently.

\begin{itemize}
\item \textbf{Above the line}, a language front-end reads source and
  \emph{plants} instructions for an abstract stack machine --- the Poplog
  VM. It does this by calling about forty ordinary Pop-11 procedures
  (\code{sysPUSHQ}, \code{sysCALL}, \code{sysPROCEDURE}, \ldots). A front-end
  is a Pop-11 program. It contains no assembly and no knowledge of any CPU.
\item \textbf{Below the line}, a back-end translates planted VM code into
  real instructions for the host architecture. This is where a port lives:
  the x86-64, AArch64, RISC-V and ARM32 code generators are each one
  directory under \code{pop/src/syscomp/}.
\end{itemize}

The two halves never meet. A language hosted on the VM inherits every port
for free --- the Forth subsystem described in Chapter~5 was written once, in
Pop-11, and runs natively on Apple Silicon, on a Raspberry Pi 5 and on a
StarFive RISC-V board with no per-platform work at all. Conversely, a new
port makes every hosted language native on that machine the day it lands.

\section{Incremental compilation, measured}

``Incremental'' here means what it says: the compiler is a procedure in the
running system, and compiling a definition is an ordinary call. There is no
build step, no link step, and no separate compile-time world.

\begin{pop}
vars t0 = systime(), k;
for k from 1 to 10000 do
    pop11_compile(stringin('define gen_' sys_>< k sys_><
                           '(x); x + ' sys_>< k sys_>< ' enddefine;'));
endfor;
npr('10000 definitions compiled in ' sys_>< ((systime() - t0) * 10) sys_>< ' ms CPU');
gen_9999(1) =>
\end{pop}
\begin{repl}
10000 definitions compiled in 150 ms CPU
** 10000
\end{repl}

Fifteen microseconds per definition, each one becoming a real, callable,
garbage-collected, \emph{natively compiled} procedure. This is the number
that explains most of what follows: when compilation is that cheap, you stop
treating it as an event. Redefining a procedure in a live session is not a
restart, it is a keystroke.

For calibration, the project's benchmark suite compares the same workloads
against CPython on identical hardware (see \code{BENCHMARKS.md} in the source
tree for methodology and the full matrix):

\begin{center}
\footnotesize
\begin{tabular}{lrrr}
\toprule
\textbf{Workload} & \textbf{Poplog (x86-64)} & \textbf{Python 3.13} & \textbf{Python 3.10}\\
\midrule
\code{nfib29} --- 1.35M recursive calls & 14.5\,ms & 3.67$\times$ slower & 8.56$\times$\\
\code{intloop10M} --- tight integer loop & 61.5\,ms & 5.60$\times$ & 6.26$\times$\\
\code{lists} --- 1M cons cells, allocate + walk & 11.4\,ms & 7.20$\times$ & 7.76$\times$\\
\code{closures1M} --- 100k closures, 1M calls & 22.0\,ms & 2.13$\times$ & 3.00$\times$\\
\code{compile500} --- runtime-compile $\times$500 & 66.7\,ms & 0.12$\times$ & 0.08$\times$\\
\bottomrule
\end{tabular}
\end{center}

The last row is the honest one, and it is worth reading carefully: Python
wins it, because \code{compile()} produces bytecode while Poplog produces
machine code. Poplog is doing strictly more work and still lands within an
order of magnitude. On the other rows --- calls, loops, allocation, closures
--- the native code shows.

\section{The saved image}

Because the VM owns the whole heap, the whole heap can be written to a file
and read back. \code{syssave} does exactly that; the resulting
\code{.psv} image contains your data \emph{and} your compiled procedures.

\begin{repl}
$ popsession checkpoint session.psv          # 10,000 procedures live in the heap
  2.9M session.psv
$ popsession restore  session.psv            # a fresh process, 64 ms including spawn
$ gen_9999(1) =>
** 10000
\end{repl}

This is how the shipped languages arrive: \code{clisp.psv},
\code{prolog.psv} and \code{pml.psv} are each an image of a Pop-11 system
that has already compiled a language on top of itself. It is also how a long-running
agent or editor session survives a restart, and why Poplog starts fast --- no
front-end is ever re-read from source at startup.

\begin{callout}{What is \emph{not} in the image}
Heap state restores; operating-system state does not. Open files, sockets,
and pointers into dynamically loaded C libraries are stale after a restore.
The discipline is simply to reopen them --- Chapter~7 returns to this.
\end{callout}

\section{Where it runs}

\begin{center}
\footnotesize
\begin{tabular}{llp{6.6cm}}
\toprule
\textbf{OS} & \textbf{Architecture} & \textbf{Status}\\
\midrule
Linux   & x86-64          & Supported --- reference platform\\
Linux   & AArch64         & Supported --- Raspberry Pi 5, MediaTek Genio 720\\
macOS   & Apple Silicon   & Supported --- native Mach-O, C\,$\leftrightarrow$\,Pop callbacks\\
Linux   & RISC-V RV64GC   & Supported --- self-hosting on StarFive VisionFive\\
Linux   & ARM32 (v6/v7)   & Supported --- long-standing port, Raspberry Pi 1--3\\
Solaris / FreeBSD & x86 / x86-64 & Supported --- upstream ports\\
Windows & x86-64          & Not yet ported (WSL2 runs the Linux build)\\
\bottomrule
\end{tabular}
\end{center}

``Supported'' means it builds, self-hosts, and runs all the languages. The
first four rows are benchmarked in this fork.
